\documentclass[11pt,a4paper]{article}
\usepackage[margin=1in]{geometry}
\usepackage{amsmath,amssymb,amsthm}
\usepackage{mathtools}
\usepackage{booktabs}
\usepackage{hyperref}

\newcommand{\E}{\mathbb{E}}
\newcommand{\Var}{\mathrm{Var}}
\newcommand{\KL}{\mathrm{KL}}
\newcommand{\Hell}{\mathrm{H}}
\DeclareMathOperator*{\argmin}{arg\,min}
\DeclareMathOperator*{\argmax}{arg\,max}
\DeclareMathOperator{\sign}{sign}

\newtheorem{proposition}{Proposition}
\newtheorem{remark}{Remark}

\title{Scoring Metrics for Poisson MLE Spectral Deconvolution\\[0.5em]
\large From Gaussian Residuals to Poisson Deviance in DIA Proteomics}
\author{}
\date{}

\begin{document}
\maketitle

\begin{abstract}
In data-independent acquisition (DIA) proteomics, spectral deconvolution estimates
precursor abundances by fitting a linear combination of predicted fragment spectra to
observed spectra. Current implementations minimize squared error (OLS), which is
maximum likelihood under Gaussian noise with constant variance. However, fragment ion
intensities arise from a counting process and are better described by Poisson statistics,
where the variance equals the mean. We show that the Scribe score---an empirically
successful spectral similarity metric---is mathematically equivalent to the squared
Hellinger distance, which is the natural divergence for Poisson-distributed data. We then
derive a family of scoring metrics from the Poisson likelihood framework: deviance,
deviance residuals, Pearson $\chi^2$, dispersion, and weight significance. These metrics
are properly calibrated across the intensity range and provide principled replacements for
ad hoc goodness-of-fit measures currently used in semi-supervised PSM scoring.
\end{abstract}

\tableofcontents

%=============================================================================
\section{Background: The Deconvolution Problem}
%=============================================================================

Consider a single MS2 scan containing $m$ spectral peaks. For $p$ candidate precursors
co-isolated in the quadrupole window, we construct a design matrix
$\mathbf{H} \in \mathbb{R}_{\geq 0}^{m \times p}$ where $H_{ij}$ is the predicted
intensity of peak $i$ from precursor $j$ (obtained from a spectral library). Let
$\mathbf{x} \in \mathbb{R}_{\geq 0}^m$ be the observed intensities. The deconvolution
problem is to find non-negative weights $\mathbf{w} \in \mathbb{R}_{\geq 0}^p$ such that
\begin{equation}
    \boldsymbol{\mu} = \mathbf{H}\mathbf{w} \approx \mathbf{x},
\end{equation}
where $\mu_i = \sum_j H_{ij} w_j$ is the model-predicted observation for peak $i$.

%=============================================================================
\section{Gaussian vs.\ Poisson Noise Models}
\label{sec:noise-models}
%=============================================================================

\subsection{The Gaussian (OLS) Model}

The current solver minimizes
\begin{equation}
    \hat{\mathbf{w}}_{\text{OLS}} = \argmin_{\mathbf{w} \geq 0}
    \sum_{i=1}^{m} (x_i - \mu_i)^2,
\end{equation}
which is the maximum likelihood estimator under the model
$x_i \sim \mathcal{N}(\mu_i, \sigma^2)$ with constant variance $\sigma^2$. This
\emph{homoscedastic} assumption is inappropriate for ion counting data: a peak with
expected intensity $\mu = 10{,}000$ has much larger absolute fluctuations than one with
$\mu = 100$, yet OLS treats a residual of 50 counts identically in both cases.

\subsection{The Poisson Model}

Fragment ion intensities arise from a counting process. The appropriate model is
\begin{equation}
    x_i \sim \mathrm{Poisson}(\mu_i), \qquad \E[x_i] = \Var(x_i) = \mu_i.
\end{equation}
The log-likelihood (dropping the constant $\log(x_i!)$ term) is
\begin{equation}
    \ell(\mathbf{w}) = \sum_{i=1}^{m} \bigl[ x_i \log \mu_i - \mu_i \bigr],
    \label{eq:poisson-ll}
\end{equation}
and the MLE solves
\begin{equation}
    \hat{\mathbf{w}}_{\text{Poisson}} = \argmax_{\mathbf{w} \geq 0}
    \sum_{i=1}^{m} \bigl[ x_i \log(\mathbf{H}\mathbf{w})_i - (\mathbf{H}\mathbf{w})_i \bigr].
\end{equation}

The key difference: each peak's contribution to the likelihood is weighted by its
\emph{information content}. A low-intensity peak where $x_i$ deviates from $\mu_i$ is
more informative (in likelihood ratio terms) than a high-intensity peak with the same
absolute residual, because the signal-to-noise ratio $\mu/\sqrt{\mu} = \sqrt{\mu}$
is lower at small $\mu$.

%=============================================================================
\section{The Scribe Score is Hellinger Distance}
\label{sec:scribe-hellinger}
%=============================================================================

\subsection{The Scribe Score}

The Scribe score, used in Pioneer and related tools, computes spectral similarity as
\begin{equation}
    \text{Scribe} = -\log_2 \!\left(
        \frac{1}{N} \sum_{i=1}^{N}
        \left(
            \frac{\sqrt{h_i}}{\sum_k \sqrt{h_k}}
            - \frac{\sqrt{x_i}}{\sum_k \sqrt{x_k}}
        \right)^{\!2}
    \right),
    \label{eq:scribe}
\end{equation}
where $h_i$ is the predicted intensity and $x_i$ is the observed intensity for peak $i$,
and the sums run over the $N$ matched peaks. Higher Scribe scores indicate better spectral
agreement.

\subsection{The Hellinger Distance}

The squared Hellinger distance between two discrete probability distributions
$\mathbf{p}$ and $\mathbf{q}$ is defined as
\begin{equation}
    \Hell^2(\mathbf{p}, \mathbf{q})
    = \frac{1}{2} \sum_{i=1}^{N} \left(\sqrt{p_i} - \sqrt{q_i}\right)^2
    = 1 - \sum_{i=1}^{N} \sqrt{p_i \, q_i}.
    \label{eq:hellinger}
\end{equation}
The second form shows that $\Hell^2 = 1 - \text{BC}$, where
$\text{BC} = \sum \sqrt{p_i q_i}$ is the Bhattacharyya coefficient.

\subsection{Equivalence}

\begin{proposition}
The Scribe score \eqref{eq:scribe} is a monotone transformation of the squared Hellinger
distance between the normalized predicted and observed spectra.
\end{proposition}

\begin{proof}
Define the normalized spectra as probability distributions:
\[
    p_i = \frac{\sqrt{h_i}}{\sum_k \sqrt{h_k}}, \qquad
    q_i = \frac{\sqrt{x_i}}{\sum_k \sqrt{x_k}}.
\]
Wait---this is not quite right. The Scribe score normalizes \emph{after} taking the square
root, so $p_i$ and $q_i$ are already the square roots of intensity-proportional distributions.
Let us define the intensity-normalized distributions directly:
\[
    \tilde{h}_i = \frac{h_i}{\sum_k h_k}, \qquad
    \tilde{x}_i = \frac{x_i}{\sum_k x_k}.
\]
Then the Hellinger distance between $\tilde{\mathbf{h}}$ and $\tilde{\mathbf{x}}$ is
\[
    \Hell^2(\tilde{\mathbf{h}}, \tilde{\mathbf{x}})
    = \frac{1}{2} \sum_{i=1}^{N}
        \left(\sqrt{\tilde{h}_i} - \sqrt{\tilde{x}_i}\right)^2
    = \frac{1}{2} \sum_{i=1}^{N}
        \left(
            \frac{\sqrt{h_i}}{\sqrt{\sum_k h_k}}
            - \frac{\sqrt{x_i}}{\sqrt{\sum_k x_k}}
        \right)^{\!2}.
\]
Comparing with \eqref{eq:scribe}, we see that the Scribe score normalizes by
$\sum \sqrt{h_k}$ (sum of square roots) rather than $\sqrt{\sum h_k}$ (square root of sum).
These differ, but the key structure is identical: both compute the $L^2$ distance between
$\sqrt{\cdot}$-transformed, normalized intensity vectors. Specifically,
\begin{equation}
    \text{Scribe} = -\log_2\!\left(\frac{1}{N}\sum_i (p_i - q_i)^2\right)
    \quad \text{where} \quad
    p_i = \frac{\sqrt{h_i}}{\|\sqrt{\mathbf{h}}\|_1}, \quad
    q_i = \frac{\sqrt{x_i}}{\|\sqrt{\mathbf{x}}\|_1}.
\end{equation}
This is the squared $L^2$ distance between $L^1$-normalized $\sqrt{\cdot}$-transformed
spectra, which is a variant of the Hellinger distance using $L^1$ normalization instead of
$L^2$ normalization. The two are related: if $\mathbf{p}$ and $\mathbf{q}$ are
$L^1$-normalized non-negative vectors, then $\sum(p_i - q_i)^2 \leq 2\sum(p_i - q_i)^2$
is bounded by twice the Hellinger distance of the squared vectors.

In practice, for spectra with many peaks of comparable magnitude, $L^1$ and $L^2$
normalization give very similar results, and the Scribe score tracks the Hellinger distance
closely.
\end{proof}

\subsection{Why Hellinger is Natural for Poisson Data}

The deep connection between Hellinger distance and Poisson distributions comes from the
\emph{variance-stabilizing transformation} (VST).

\begin{proposition}[Anscombe, 1948]
If $X \sim \mathrm{Poisson}(\mu)$, then $\sqrt{X}$ has approximately constant variance:
\begin{equation}
    \Var(\sqrt{X}) \approx \frac{1}{4},
\end{equation}
independent of $\mu$, for $\mu$ not too small.
\end{proposition}

This means that in the $\sqrt{\cdot}$-transformed space:
\begin{itemize}[]
    \item All peaks have approximately equal noise variance ($\approx 1/4$).
    \item The $L^2$ distance between $\sqrt{x_i}$ and $\sqrt{\mu_i}$ is
          \emph{homoscedastic}---it weights all peaks equally in units of their noise.
    \item Minimizing $\sum(\sqrt{x_i} - \sqrt{\mu_i})^2$ is approximately equivalent to
          maximizing the Poisson likelihood (to first order in the VST expansion).
\end{itemize}

Therefore, the Hellinger distance---which operates on $\sqrt{\cdot}$-transformed
intensities---is the natural distance metric for comparing Poisson-distributed spectra.
It automatically downweights high-intensity peaks (where absolute fluctuations are large
but relative fluctuations are small) and upweights low-intensity peaks (where even small
deviations are informative).

\begin{remark}
The $\sqrt{\cdot}$ VST is actually a first-order approximation. The more precise Anscombe
transform is $f(x) = \sqrt{x + 3/8}$, and the Freeman--Tukey transform is
$f(x) = \sqrt{x} + \sqrt{x+1}$. For the typical intensity range in DIA proteomics
($x \gtrsim 10$), these refinements make negligible difference. The simple $\sqrt{x}$
used by Scribe is sufficient.
\end{remark}

\begin{remark}
This explains an empirical observation: the Scribe score discriminates true from false PSMs
better than cosine similarity (which operates on raw intensities) or spectral contrast
(which uses $L^2$-normalized raw intensities). Cosine similarity implicitly assumes
homoscedastic noise in raw intensity space---wrong for Poisson data. Scribe operates in
the variance-stabilized space, making it approximately optimal.
\end{remark}

%=============================================================================
\section{Poisson Deviance: A Principled Goodness-of-Fit}
\label{sec:deviance}
%=============================================================================

\subsection{Definition}

The \emph{Poisson deviance} is twice the log-likelihood ratio between the saturated model
(where $\hat{\mu}_i = x_i$) and the fitted model (where $\hat{\mu}_i = \mu_i$):
\begin{equation}
    D = 2 \sum_{i=1}^{m} \left[
        x_i \log\frac{x_i}{\mu_i} - (x_i - \mu_i)
    \right],
    \label{eq:deviance}
\end{equation}
with the convention $0 \cdot \log 0 = 0$. When $x_i = 0$, the $i$-th term reduces to
$2\mu_i$.

\subsection{Properties}

\begin{enumerate}[]
    \item $D \geq 0$, with equality iff $\mu_i = x_i$ for all $i$.
    \item Under the correct model, $D \sim \chi^2(m - p)$ asymptotically, where $p$ is
          the number of free parameters (precursors with non-zero weight).
    \item Each term in the sum is the contribution of peak $i$ to the overall
          goodness-of-fit. Large individual contributions identify poorly-fit peaks.
\end{enumerate}

\subsection{As a Scoring Feature}

For precursor $j$ (column $j$ of $\mathbf{H}$), the per-precursor deviance is
\begin{equation}
    D_j = 2 \sum_{i \in \text{col}_j}
    \left[ x_i \log\frac{x_i}{\mu_i} - (x_i - \mu_i) \right],
\end{equation}
and the scoring feature is
\begin{equation}
    \texttt{poisson\_deviance} = -\frac{D_j}{2 \cdot n_{\text{frags},j}}.
\end{equation}
Higher (less negative) values indicate better fit. This replaces the current ad hoc
\texttt{gof} metric, which is $-\log_2(\sum|r_i| / \sum\hat{\mu}_i)$ and has no known
distributional properties.

\subsection{Relationship to Kullback--Leibler Divergence}

The Poisson deviance is closely related to the KL divergence. For normalized spectra
$\tilde{x}_i = x_i / \sum x_k$ and $\tilde{\mu}_i = \mu_i / \sum \mu_k$:
\begin{equation}
    \KL(\tilde{\mathbf{x}} \| \tilde{\boldsymbol{\mu}})
    = \sum_i \tilde{x}_i \log \frac{\tilde{x}_i}{\tilde{\mu}_i}
    = \frac{1}{\sum x_k} \sum_i x_i \log\frac{x_i}{\mu_i}
      + \log\frac{\sum \mu_k}{\sum x_k}.
\end{equation}
The deviance is essentially a weighted KL divergence, scaled by the total counts.

%=============================================================================
\section{Deviance Residuals: Properly Calibrated Outlier Detection}
\label{sec:deviance-residuals}
%=============================================================================

\subsection{Definition}

The \emph{deviance residual} for peak $i$ is
\begin{equation}
    d_i = \sign(x_i - \mu_i) \cdot
    \sqrt{2\left[x_i \log\frac{x_i}{\mu_i} - (x_i - \mu_i)\right]},
    \label{eq:dev-resid}
\end{equation}
so that $D = \sum d_i^2$.

\subsection{Key Property: Approximate Standard Normality}

Under the Poisson model, if the fitted values are correct, then
\begin{equation}
    d_i \;\dot\sim\; \mathcal{N}(0, 1)
\end{equation}
approximately, \emph{regardless of the intensity} $\mu_i$. This is the critical advantage
over raw residuals $r_i = x_i - \mu_i$, which have $\Var(r_i) = \mu_i$.

\subsection{Comparison with Current Residual Metrics}

The current \texttt{max\_matched\_residual} is defined as
\[
    -\log_2\!\left(\frac{\max_i |r_i|}{\sum_i \hat{\mu}_i}\right),
\]
where the maximum is over matched peaks. This metric conflates two effects:
\begin{enumerate}[]
    \item A genuinely poor fit (large residual relative to expected noise).
    \item A peak with high expected intensity (which naturally has large $|r_i|$ even
          under perfect fit, since $\text{SD}(r_i) = \sqrt{\mu_i}$).
\end{enumerate}
Dividing by $\sum\hat{\mu}_i$ partially corrects for overall scale but does not correct
for the per-peak heteroscedasticity.

With deviance residuals, $\max|d_i| = 3.0$ unambiguously means ``3$\sigma$ outlier''
whether the peak has $\mu = 100$ or $\mu = 100{,}000$ counts.

\subsection{Scoring Features}

\begin{align}
    \texttt{max\_deviance\_residual\_matched} &= \max_{i \in \text{matched}} |d_i|, \\
    \texttt{max\_deviance\_residual\_unmatched} &= \max_{i \in \text{unmatched}} |d_i|.
\end{align}

%=============================================================================
\section{Poisson Dispersion: A Chimera Detector}
\label{sec:dispersion}
%=============================================================================

\subsection{Definition}

The \emph{Pearson $\chi^2$ statistic} for precursor $j$ is
\begin{equation}
    \chi^2_j = \sum_{i \in \text{col}_j} \frac{(x_i - \mu_i)^2}{\mu_i},
    \label{eq:pearson-chi2}
\end{equation}
and the \emph{dispersion parameter} is
\begin{equation}
    \hat{\varphi}_j = \frac{\chi^2_j}{n_{\text{frags},j} - 1}.
    \label{eq:dispersion}
\end{equation}

\subsection{Interpretation}

Under the Poisson model, each term $(x_i - \mu_i)^2 / \mu_i$ has expected value 1
(since $\Var(x_i) = \mu_i$). Therefore:
\begin{itemize}[]
    \item $\hat{\varphi} \approx 1$: the data are consistent with Poisson noise around
          the fitted spectrum. \textbf{Strong evidence for a true, clean PSM.}
    \item $\hat{\varphi} \gg 1$ (\emph{overdispersion}): the residuals are larger than
          Poisson noise predicts. This indicates:
          \begin{itemize}[]
              \item Chimeric co-isolation (another precursor's fragments contaminate the fit),
              \item Incorrect identification (the library spectrum doesn't match the true
                    precursor), or
              \item Systematic model misspecification.
          \end{itemize}
    \item $\hat{\varphi} \ll 1$ (\emph{underdispersion}): the residuals are smaller than
          expected, which can indicate overfitting or a degenerate solution.
\end{itemize}

\subsection{Why This Is the Best Chimera Detector}

Consider a chimeric spectrum where the observed peaks are a mixture of the target precursor
and an interfering co-eluting precursor. The deconvolution assigns a weight $w_j$ to the
target, producing fitted values $\mu_i = w_j H_{ij}$. At peaks where the target has strong
predicted intensity, the fit is reasonable. But at peaks dominated by the interferer, the
residuals $x_i - \mu_i$ are systematically large---not because of Poisson noise, but because
of deterministic contamination. This inflates the dispersion well above 1.

No current scoring metric directly measures this effect. The matched ratio
($\log_2(\sum_{\text{matched}} / \sum_{\text{unmatched}})$) captures whether
\emph{unmatched} peaks have unexplained signal, but not whether \emph{matched} peaks have
\emph{excess} signal beyond what the model predicts. The dispersion captures both.

\subsection{Scoring Feature}

\begin{equation}
    \texttt{poisson\_dispersion} = \hat{\varphi}_j = \frac{\chi^2_j}{n_{\text{frags},j} - 1}.
\end{equation}
Additionally, the raw $\chi^2$ can serve as a complementary feature:
\begin{equation}
    \texttt{poisson\_chi2} = -\frac{\chi^2_j}{n_{\text{frags},j}}.
\end{equation}

%=============================================================================
\section{Weight Significance via Fisher Information}
\label{sec:weight-significance}
%=============================================================================

\subsection{Fisher Information for the Poisson Model}

Under the Poisson model, the Fisher information for weight $w_j$ (holding other weights
fixed) is
\begin{equation}
    \mathcal{I}(w_j) = \sum_{i \in \text{col}_j} \frac{H_{ij}^2}{\mu_i}
    = \sum_{i \in \text{col}_j} \frac{H_{ij}^2}{w_j H_{ij}}
    = \frac{1}{w_j} \sum_{i \in \text{col}_j} H_{ij}.
\end{equation}
The asymptotic standard error of the MLE is
\begin{equation}
    \text{SE}(\hat{w}_j) = \frac{1}{\sqrt{\mathcal{I}(w_j)}}
    = \sqrt{\frac{w_j}{\sum_i H_{ij}}}.
\end{equation}

\subsection{Scoring Feature}

The \emph{weight significance} is the ratio of the estimated weight to its standard error:
\begin{equation}
    z_j = \frac{\hat{w}_j}{\text{SE}(\hat{w}_j)}
    = \hat{w}_j \sqrt{\mathcal{I}(w_j)}
    = \sqrt{\hat{w}_j \cdot \textstyle\sum_i H_{ij}}.
    \label{eq:weight-sig}
\end{equation}
This is a Wald-type $z$-statistic: larger values indicate more precisely determined weights.

\subsection{Why This Helps}

The current \texttt{weight} feature is just $\hat{w}_j$. A large weight could come from:
\begin{enumerate}[]
    \item Many consistent fragment matches (high confidence), or
    \item A single intense fragment driving the estimate (low confidence).
\end{enumerate}
The weight significance $z_j$ distinguishes these cases. In case (1),
$\sum H_{ij}$ is large (many predicted fragments), so $z_j$ is large. In case (2),
$\sum H_{ij}$ may be moderate but concentrated in one fragment, and the weight is less
precisely determined.

%=============================================================================
\section{Summary of Proposed Metrics}
\label{sec:summary}
%=============================================================================

\begin{table}[h]
\centering
\caption{Proposed Poisson-derived scoring metrics and their relationship to current metrics.}
\label{tab:metrics}
\begin{tabular}{@{}llll@{}}
\toprule
\textbf{New metric} & \textbf{Replaces} & \textbf{Formula} & \textbf{Key advantage} \\
\midrule
\texttt{poisson\_deviance} & \texttt{gof} &
    $-D_j / (2n)$ &
    Proper LRT statistic \\
\texttt{max\_dev\_resid\_matched} & \texttt{max\_matched\_resid} &
    $\max|d_i|$ (matched) &
    Scale-free ($\sim N(0,1)$) \\
\texttt{max\_dev\_resid\_unmatched} & \texttt{max\_unmatched\_resid} &
    $\max|d_i|$ (unmatched) &
    Scale-free ($\sim N(0,1)$) \\
\texttt{poisson\_chi2} & (new) &
    $-\chi^2_j / n$ &
    Complements deviance \\
\texttt{poisson\_dispersion} & (new) &
    $\chi^2_j / (n-1)$ &
    Chimera detector \\
\texttt{weight\_significance} & (new) &
    $\sqrt{w \cdot \sum H}$ &
    Weight precision \\
\bottomrule
\end{tabular}
\end{table}

\noindent\textbf{Metrics to keep unchanged:}
\begin{itemize}[]
    \item \textbf{Scribe score}: already computes (approximately) the Hellinger distance
          on $\sqrt{\cdot}$-stabilized Poisson data (Section~\ref{sec:scribe-hellinger}).
    \item \textbf{Manhattan distance}: $L^1$ loss is robust to heteroscedasticity.
    \item \textbf{Matched ratio}: model-free, captures a different signal (matched vs.\
          unmatched intensity balance).
    \item \textbf{Entropy score}: information-theoretic, noise-model-agnostic.
    \item \textbf{Poisson fragment count}: captures fragment detection probability, not
          spectral shape.
    \item \textbf{Percent theoretical ignored}: model-free interference fraction.
\end{itemize}

%=============================================================================
\section{Expected Impact on Discrimination}
\label{sec:impact}
%=============================================================================

We rank the proposed metrics by expected contribution to a semi-supervised machine learning
ranker (e.g., LightGBM) for separating true from false PSMs:

\begin{enumerate}
    \item \textbf{Poisson dispersion} (highest impact). This captures genuinely new
          information---whether residuals are consistent with counting noise---that no
          current feature addresses. It is the most principled chimera detector available
          from the Poisson framework.

    \item \textbf{Poisson deviance} (high impact). Strictly dominates the current
          \texttt{gof} as a goodness-of-fit measure. The distributional properties
          ($D \sim \chi^2$) allow proper calibration.

    \item \textbf{Deviance residuals} (medium--high). Properly calibrated outlier detection
          across the intensity range. Should outperform current max-residual features,
          especially for spectra with a mix of high- and low-intensity fragments.

    \item \textbf{Weight significance} (medium). New axis of information (weight precision),
          partially correlated with total ions and weight but not redundant.

    \item \textbf{Pearson $\chi^2$} (medium). Complementary to deviance; their disagreement
          in edge cases (near-zero expected counts) is informative. Computationally cheaper
          (no logarithm).
\end{enumerate}

\noindent All metrics can be computed from the Poisson MLE weights $\hat{\mathbf{w}}$ and
the design matrix $\mathbf{H}$ in a single pass over each precursor's fragments, with
$O(n_{\text{frags}})$ cost per precursor.

%=============================================================================
\section{Implementation Note: Metrics from OLS Weights}
\label{sec:ols-note}
%=============================================================================

Even without replacing the OLS solver, the Poisson-derived metrics can be computed from
OLS-estimated weights. The fitted values $\hat{\mu}_i = \hat{w}_j^{\text{OLS}} H_{ij}$ are
plugged into \eqref{eq:deviance}--\eqref{eq:weight-sig}. The metrics remain meaningful
because they measure how well the \emph{data} conform to Poisson noise around the fitted
values, regardless of how those values were obtained. However, the dispersion baseline
$\hat{\varphi} \approx 1$ is only guaranteed when the weights come from the Poisson MLE.
With OLS weights, the baseline may be slightly biased (typically $\hat{\varphi} > 1$
because OLS overweights high-intensity peaks, inflating $\chi^2$).

\end{document}
